%% =====================================================================
%%  B4-T-16-01 -- what B4 contributes to "Concession and Retreat"
%%  -------------------------------------------------------------------
%%  LAYER A: APPEND-ONLY. Written once, then left alone. Material found
%%  only here sits in the `unique` slot and is protected by rule R10.
%% =====================================================================
%% @kind: source
%% @id: B4-T-16-01
%% @book: B4
%% @topic: T-16-01
%% @locator: ch. 2, pp. 43--86
%% @status: distilled
%% @unique: yes
%% @integrated: yes
%% @updated: 2026-09-19

\dossier{B4}{T-16-01}{\bookshort{B4}}{ch. 2, pp. 43--86}

\begin{distilled}
  \item A large request followed by a smaller one produces more compliance with the smaller than the smaller request produces alone.
  \item The effect is attributed to two mechanisms acting together: the reciprocal obligation created by the other party's apparent concession, and the contrast that makes the second request look moderate against the first.
  \item The technique also improves how the resulting agreement is felt about afterwards, and raises the rate at which it is actually honoured.
\end{distilled}

\begin{unique}
  \item The reciprocal-concession mechanism as distinct from contrast alone, and the observation that both fire on the form of the move rather than on the intent behind it.
  \item The finding that compliance obtained this way is more durable and more likely to be followed through than compliance obtained by the smaller request directly.
\end{unique}
